\documentclass[preview]{standalone}

\usepackage{amsmath}
\usepackage{amssymb}
\usepackage{stellar}
\usepackage{definitions}

\begin{document}

\id{psicologia-espressione-emotiva}
\genpage

\section{Espressione e contesto emotivo}

\begin{snippet}{darwin-espressioni-emozioni}
    Charles Darwin, in \emph{L'espressione delle emozioni nell'uomo e negli
    animali}, si chiese se le espressioni facciali delle emozioni fossero
    culturalmente invarianti. Nella sua prospettiva, le emozioni sono stati
    mentali ereditati e specializzati, formatisi per rispondere a classi
    ricorrenti di situazioni.
\end{snippet}

\begin{snippetdefinition}{programma-affettivo-definition}{Programma affettivo}
    Un \emph{programma affettivo} è uno specifico programma nervoso, evolutosi
    nel tempo, che regola un'emozione di base e consente alla specie di
    adattarsi efficacemente al proprio ambiente.
\end{snippetdefinition}

\begin{snippet}{programmi-affettivi}
    Secondo Tomkins, le emozioni sono associate a scopi universali connessi alla
    sopravvivenza dell'individuo e della specie. Questa prospettiva, sostenuta
    anche da Ekman, Izard e Panksepp, afferma che ogni emozione di base sia
    regolata da uno specifico programma affettivo.
\end{snippet}

\begin{snippetdefinition}{ipotesi-universalita-espressioni-definition}{Ipotesi dell'universalità delle espressioni emotive}
    L'\emph{ipotesi dell'universalità delle espressioni emotive} sostiene che
    alcune espressioni facciali delle emozioni di base siano innate, universali
    e riconoscibili in tutte le culture.
\end{snippetdefinition}

\begin{snippet}{emozioni-base-ekman}
    Secondo Ekman, le emozioni di base includono collera, disgusto, paura,
    gioia, tristezza e sorpresa. L'ipotesi dell'universalità riguarda queste
    emozioni di base, non tutte le espressioni facciali né ogni modalità
    culturale di manifestazione emotiva.
\end{snippet}

\begin{snippetnote}{universalita-emozioni-base-note}{Universalità ed emozioni di base}
    L'ipotesi dell'universalità riguarda soltanto le emozioni di base. Ekman non
    sostiene che tutte le espressioni facciali siano universali o che tutte le
    culture esprimano ogni emozione nello stesso modo.
\end{snippetnote}

\includesnpt[width=62\%,src=/snippet/static-psychology/emotional-expressions-context.jpg]{centered-img}

\begin{snippetdefinition}{teoria-neuroculturale-definition}{Teoria neuroculturale}
    La \emph{teoria neuroculturale} sostiene che l'espressione delle emozioni
    dipenda dal contributo congiunto del cervello, modellato dall'evoluzione, e
    della cultura.
\end{snippetdefinition}

\begin{snippet}{biologia-cultura-espressione}
    L'ipotesi di un legame tra emozioni ed espressioni facciali non esclude
    variazioni culturali significative. È quindi più appropriato parlare di
    interdipendenza tra componenti biologiche e influenze culturali nella
    manifestazione delle emozioni.
\end{snippet}

\begin{snippetdefinition}{prospettiva-contestualista-definition}{Prospettiva contestualista}
    La \emph{prospettiva contestualista} sostiene che le espressioni facciali
    assumano un valore emotivo definito soltanto in riferimento a una situazione
    specifica.
\end{snippetdefinition}

\begin{snippetdefinition}{indessicalita-definition}{Indessicalità}
    L'\emph{indessicalità} è la proprietà per cui un segno, come un'espressione
    facciale, rimanda a una realtà specifica attraverso l'uso sistematico di
    indizi contestuali.
\end{snippetdefinition}

\begin{snippet}{manifestazione-emotiva-multimodale}
    Le espressioni facciali sono normalmente accompagnate da movimenti corporei,
    gesti e informazioni vocali. La percezione delle manifestazioni emotive è
    quindi multimodale e dipende anche dal contesto fisico e sociale, che può
    facilitare o modificare il riconoscimento delle emozioni.
\end{snippet}

\end{document}
